\section{Types of sprinkler systems}
In \kw{sprinkler irrigation} the water is carried under pressure in pipes and sprayed through nozzles, so that it falls on the crop like rain. The systems are classified by their mobility.
\begin{exlist}
\item \kw{Hand-move lateral.} Aluminium pipes with couplers are moved by hand after each set. Low capital cost, high labour.
\item \kw{Solid-set.} The laterals and sprinklers stay in the field for the whole season (or are buried), and are controlled by valves; also used for frost protection and cooling.
\item \kw{Side-roll (wheel-line).} A lateral on large wheels is rolled to the next position; for low crops on rectangular fields.
\item \kw{Centre pivot.} A lateral on towers turns round a fixed centre, and irrigates a circle; one of the most widely used systems for large fields. The speed of the outer tower is greatest, so the nozzles are larger towards the end, with a discharge proportional to the radius.
\item \kw{Linear move.} The lateral moves in a straight line, and irrigates a rectangle.
\item \kw{Travelling gun (big gun).} One large sprinkler on a hose reel; flexible, but with high pressure and large drops.
\end{exlist}

\section{Components of a sprinkler system}
A pump and a power unit; the \emph{mainline} and the \emph{sub-main}, which can be buried; the \emph{lateral}, which carries the sprinklers; the \emph{riser}, which lifts the sprinkler above the crop; and the \emph{sprinkler head}, with one or two \emph{nozzles}. The head turns by the impact of a spring-loaded arm (\emph{impact} sprinkler) or by a gear drive. Small nozzles, at 200 to 300 kPa, give fine drops; large nozzles for big guns need 400 to 700 kPa.

\section{Advantages and disadvantages of sprinkler irrigation}
\textbf{Advantages:} it can be used on any slope without levelling; the application rate can be matched to the infiltration of a light soil; it gives an even distribution if the design is right; it can be automated; the depth is easy to control, and small depths can be given for germination, cooling and frost protection; it saves land (no channels).
\par\noindent \textbf{Disadvantages:} the energy for the pressure; a higher capital cost than surface irrigation; the wind distorts the pattern, and the drops can be blown away (drift loss); evaporation in the air; the wetting of leaves can promote disease and can scorch the leaves if the water is saline; and the large drops of a big gun can seal the surface of a bare soil.

\section{Drops, energy and the soil surface}
The drops that leave the nozzle have a diameter between 0.5 and 4 mm. A large drop falls faster and has more \kw{kinetic energy}; on a bare soil it breaks the aggregates, seals the surface and lowers the infiltration. A fine drop is blown by the wind and loses more by evaporation. The nozzle pressure, the nozzle size and the type of head control the spectrum of drops: a high pressure with a small nozzle gives fine drops, a low pressure with a large nozzle gives coarse ones. The aim is the largest drop that does not damage the soil and the crop.

\section{Wind and the pattern of a sprinkler}
Without wind a sprinkler gives a pattern that is roughly triangular in section, with the most water near the head, falling to nothing at the radius. Wind shifts the pattern downwind and shortens it upwind, and a wind above 3 m\,s$^{-1}$ requires closer spacing along the wind direction. The recommended spacing, as a fraction of the wetted diameter, is 65 per cent for a wind below 2 m\,s$^{-1}$, 60 per cent for 2 to 4 m\,s$^{-1}$ and 50 per cent above that. Overlapping patterns add up to an even depth, which the uniformity coefficient measures (Section~\ref{ch:cu}).

\section{Application rate and the infiltration of the soil}
The rate at which the sprinklers apply water must not be larger than the soil can take, or water will pond and run off. The soil takes water at a rate that falls to a final value, and the maximum allowable rate depends on the soil and the slope; typical limits are 20 to 25 mm\,h$^{-1}$ in a sand, 10 to 15 in a loam and 5 or less in a clay, reduced by 20 to 40 per cent on slopes. If the rate is too high, a smaller nozzle, a wider spacing or a different sprinkler must be chosen; the check is the last of the ten tests of this chapter.
